\documentclass{article}
\usepackage[utf8]{inputenc}
\usepackage[top=2cm, bottom=2cm, left=2cm, right=2cm]{geometry}
\usepackage[hidelinks]{hyperref}
\usepackage[showans=true, showkey=true]{pro_book}

\usepackage{polyglossia}
\setmainlanguage{vietnamese}
\usepackage{fontspec}
\setmainfont{Times New Roman}

\title{Demo Sách Chuyên Nghiệp với pro\_book.sty}
\author{Antigravity}
\date{\today}

\begin{document}

\maketitle
\tableofcontents
\StartProBook

\section{Lý Thuyết (Premium Design)}

\begin{dn}{Hàm số lũy thừa}{def:luythua}
    Hàm số $y=x^\alpha$ với $\alpha \in \mathbb{R}$ được gọi là hàm số lũy thừa.
\end{dn}

\begin{dl}{Đạo hàm}{thm:daoham}
    $(x^\alpha)' = \alpha x^{\alpha-1}$.
\end{dl}

\begin{note}
    Chú ý: Công thức trên chỉ đúng khi $x > 0$.
\end{note}

\section{Ví Dụ Minh Họa}

\begin{vd}{Tìm tập xác định}{vd:txd}
    Tìm tập xác định của hàm số $y=(x-1)^{\frac{1}{3}}$.
    \loigiai{
        Vì số mũ $\frac{1}{3}$ không nguyên nên điều kiện là $x-1 > 0 \Leftrightarrow x > 1$.
        Vậy $D = (1; +\infty)$.
    }
\end{vd}

\section{Bài Tập Trắc Nghiệm}

\begin{bt}{Tính đạo hàm}{bt:daoham}
    Tính đạo hàm của $y=x^{\pi}$.
    \choice
    {$\pi x^{\pi-1}$}
    {\True $\pi x^{\pi-1}$}
    {$x^{\pi} \ln x$}
    {$\pi x^{\pi}$}
    \loigiai{
        Áp dụng công thức $(x^\alpha)' = \alpha x^{\alpha-1}$, ta có $y' = \pi x^{\pi-1}$.
    }
\end{bt}

\begin{ex}[Câu hỏi nhận biết]
    Đồ thị hàm số $y=\frac{x-1}{x+1}$ có tiệm cận đứng là:
    \choice
    {$x=1$}
    {\True $x=-1$}
    {$y=1$}
    {$y=-1$}
    \loigiai{
        Tiệm cận đứng là nghiệm của mẫu số: $x+1=0 \Rightarrow x=-1$.
    }
\end{ex}

\begin{ex}[Câu hỏi vận dụng]
    Cho hàm số $y=f(x)$ có bảng biến thiên như sau...
    Hàm số đồng biến trên khoảng nào?
    \choice
    {$(-\infty; -1)$}
    {$(0; 1)$}
    {\True $(1; +\infty)$}
    {$(-1; 1)$}
    \loigiai{
        Dựa vào bảng biến thiên, $y' > 0$ trên $(1; +\infty)$.
    }
\end{ex}

\section{Trắc Nghiệm Đúng Sai}

\begin{ex}[Đúng sai]
    Cho hàm số $y=x^3 - 3x^2 + 2$.
    \choiceTF
    {Hàm số có 2 cực trị.}
    {\True Hàm số đạt cực đại tại $x=0$.}
    {Giá trị cực tiểu là $-2$.}
    {\True Đồ thị cắt trục tung tại điểm có tung độ bằng 2.}
    \loigiai{
        $y' = 3x^2 - 6x$. $y'=0 \Leftrightarrow x=0$ hoặc $x=2$.
        Tại $x=0, y=2$ (Cực đại).
        Tại $x=2, y=-2$ (Cực tiểu).
    }
\end{ex}

\StopProBook
\PrintAnswers

\end{document}
